\documentclass[../../main.tex]{subfiles}
\begin{document}

\begin{center}
  {\fontsize{48pt}{58pt}\selectfont \textbf{1973}\normalsize\textbf{年度}}
\end{center}

\begin{table}[htb]
  \centering
  \caption{出題テーマと難易度}
  \begin{tabular}{|c|p{9cm}|c|}
    \hline
    \textbf{問題番号} & \textbf{テーマ} & \textbf{難易度} \\
    \hline
    第1問 & 球面による遮蔽と等速円運動 & \\
    \hline
    第2問 & 0，1，2からなる数列の累乗和 & \\
    \hline
    第3問 & 区分的な関数の最大値と最小値の差 & \\
    \hline
    第4問 & 重なる正方形の面積と中心の範囲 & \\
    \hline
    第5問 & 不等式で定まる領域の面積の導関数 & \\
    \hline
    第6問 & 放物線上の点における接線の傾き & \\
    \hline
  \end{tabular}
\end{table}

\begin{table}[htb]
  \centering
  \caption{解答}
  \renewcommand{\arraystretch}{1.4}
  \begin{tabular}{|c|c|p{9cm}|}
    \hline
    \textbf{問題番号} & \textbf{小問} & \textbf{答え} \\
    \hline
    第1問 & -- & $8$秒間見える．$NP$の最大値は$\sqrt{15}a$，最小値は$\sqrt{13}a$ \\
    \hline
    第2問 & -- & $f_k=(2-2^{k-1})f_1+(2^{k-1}-1)f_2$ \\
    \hline
    第3問 & (1) & 最小値は$\dfrac{1}{2}$ \\
    \cline{2-3}
    & (2) & $a=\dfrac{1}{2}$ \\
    \hline
    第4問 & -- & 点$P$の範囲は解答の図を参照．面積は$2(1-\log 2)$ \\
    \hline
    第5問 & -- & $f'(t)=\dfrac{2(t^2-1)}{(1+t^2)^2}\log t$ \\
    \hline
    第6問 & -- & 接線の傾きは$\dfrac{2v-u}{u}$ \\
    \hline
  \end{tabular}
\end{table}

\paragraph{試験形態}
6問．試験時間は要確認．

\end{document}
